
\subsubsection{N-gram Overlap Detection}

The dominant approach to contamination detection relies on exact substring matching. Brown et al. (2020) introduced 13-gram overlap in GPT-3, flagging training documents sharing 13+ consecutive tokens with benchmark items. OpenAI's GPT-4 technical report refined this to 50-character overlap thresholds. These methods scale efficiently via suffix arrays and Bloom filters, enabling analysis of trillion-token corpora.

However, n-gram methods fail on paraphrased contamination. Yang et al. (2023) demonstrate that training Llama-2-13B on rephrased MMLU achieves 85.9\% accuracy while evading n-gram detection entirely. The fundamental limitation is definitional: exact matching cannot capture semantic equivalence across phrasings.

\subsubsection{Quiz-Based and Behavioral Methods}

Golchin and Surdeanu's Data Contamination Quiz (DCQ) takes a behavioral approach: generate completion prompts from benchmark items and measure whether models reproduce memorized content with suspiciously high accuracy. DCQ achieves higher memorization detection rates than n-gram methods and works without training data access. The limitation is scalability---DCQ provides binary output per item and requires individual probing.

\subsubsection{Surveys on Contamination}

Recent surveys document contamination levels up to 45\% in popular LLMs on common benchmarks (Sainz et al., 2024; Li et al., 2024). Xu et al. (2024) evaluate detection method assumptions and limitations. These works establish the scope of the problem but do not propose paraphrase-resistant continuous metrics.

\subsubsection{This Work's Position}

This work addresses a gap: no existing method provides scalable, continuous-valued, paraphrase-resistant contamination measurement. The approach differs by measuring the consequence of diverse training exposure---confidence uniformity---rather than matching training content directly.
